\section{منطق انتخاب محدودیت‌ها (\lr{Based On What?})}

\begin{itemize}

    \item\textbf{\lr{Ready for Pull (WIP = 6)}:}
این مقدار بر اساس ظرفیت ۴ نفره تیم و اصل \lr{Just-In-Time (JIT)} انتخاب شده است. تخصیص ۱٫۵ کارت به ازای هر نفر تضمین می‌کند که تیم همواره کار شفاف و آماده برای شروع دارد، اما برنامه‌ریزی بیش از حد زودهنگام (\lr{Over-planning}) که منجر به اتلاف (\lr{Waste}) می‌شود، رخ نمی‌دهد.

    \item\textbf{\lr{Modeling(Design, Domain \& Architectural modeling) (WIP = 2)}:}
از آنجایی که هسته این نرم‌افزار بر پایه \lr{ddd} است، طراحی \lr{aggregate}ها نیاز به تفکر عمیق دارد و از طرفی محصول باید برای ذی نفع نهایی نیز طراحی مناسبی داشته باشد. محدودیت ۲ کارت از \lr{over-design} همزمان جلوگیری کرده و تضمین می‌کند تصمیمات طراحی محور دقیقاً در زمان نیاز گرفته و بلافاصله وارد فاز پیاده‌سازی می‌شوند.z
    \item\textbf{\lr{Doing (WIP = 4)}:}
این عدد دقیقاً معادل تعداد اعضای تیم است (اصل \lr{One-Piece Flow}). هر توسعه‌دهنده در هر لحظه تنها مجاز به کار روی \textbf{یک} تسک فعال است. این محدودیت، جلوی چندوظیفگی (\lr{Multitasking}) را می‌گیرد و سرعت تحویل (\lr{Lead Time}) هر ویژگی را به شدت کاهش می‌دهد.

    \item\textbf{\lr{Code Review \& Architecture Check (WIP = 2)}:}
بازبینی کد در این پروژه صرفاً بررسی سینتکس نیست. در این ستون، تمرکز بسیار ویژه‌ای بر رعایت ایزوله‌سازی دامنه و \textbf{جداسازی دقیق و سخت‌گیرانه منطق لایه \lr{Application} از \lr{Infrastructure}} وجود دارد. کدهای مربوط به دیتابیس مطلقاً نباید به لایه \lr{Application} نشت کنند. محدودیت سفت و سختِ ۲ کارت تضمین می‌کند که این بررسی‌های معماری حیاتی با بالاترین دقت انجام شوند و در صورت انباشت کارت‌ها، کل تیم توسعه را متوقف کرده و به رفع این گلوگاه ملزم می‌کند (\lr{Swarming}).

    \item\textbf{\lr{Testing / Validation (WIP = 3)}:}
با توجه به اینکه تیم تستر اختصاصی ندارد، این عدد به گونه‌ای تنظیم شده که اعضا را مجبور کند به جای شروع کدهای جدید، کارهای توسعه‌یافته را نهایی کنند.    
    

\end{itemize}
